% Chapter 5: at the proposal stage this chapter holds the work plan, risks and fallbacks.
% At final submission it becomes the conclusion: findings, contributions, limitations and
% future work.

\section{Summary of the Proposal}
This thesis proposes to detect abnormal trading days on the Dhaka Stock Exchange with a
deep autoencoder whose inputs are returns standardized by GARCH-family conditional
volatility, and to train it on global market indices before transferring it to the DSE. The
work tests three hypotheses: that standardization improves detection and reduces false
alarms in volatile periods, that detectors transfer from global markets to the DSE, and that
point adjustment inflates reported performance on financial data. Phase~0 of the project is
complete: DSE data sources have been identified, 23 papers have been reviewed, the code
repository is set up, and a first volatility model run on S\&P~500 data has shaped the
design.

\section{Work Plan}
The work proceeds in five phases, each ending with a milestone that must be reached before
the next phase depends on it. Table~\ref{tab:plan} lists them in order.

\begin{table}[H]
\caption{Work plan.}
\label{tab:plan}
\begin{center}
\begin{tabular}{p{0.5cm} p{2.6cm} p{6.0cm} p{4.0cm}}
\hline
 & Phase & Work & Milestone \\
\hline
1 & Validation & DSE data sources identified (download not yet run); literature review
of 23 papers; first GARCH run; code repository and seed policy & Scope agreed with the
supervisor \\
2 & Proposal and setup & Research questions and hypotheses; global and DSE data collection;
DSE quality check & Proposal defended; DSE data pass the go/no-go criteria, or the
fallback market is chosen \\
3 & Data and labels & Cleaning (missing days, halts, floor prices, bonus shares, price
limits); GARCH, EGARCH and GJR fits with diagnostics; known-event, residual and injected
labels & Labelling rules fixed before any detector is trained \\
4 & Models and experiments & Baselines; raw versus standardized autoencoder; transfer
from global indices to the DSE; five-seed evaluation; regime analysis & All results
reported over five seeds \\
5 & Writing and submission & Thesis book; code release; data release if licensing allows;
IEEE paper draft & Thesis submitted \\
\hline
\end{tabular}
\end{center}
\end{table}

\section{Division of Work}
Each team member owns one strand of the work:
\begin{itemize}
    \item \textbf{\StudentTwoName:} data collection, cleaning and labels.
    \item \textbf{\StudentThreeName:} GARCH-family modelling and baselines.
    \item \textbf{\StudentOneName:} deep models and transfer experiments.
\end{itemize}
All three members write the thesis, with one member responsible for final integration.

\section{Risks and Fallbacks}\label{sec:risks}
\begin{itemize}
    \item \textbf{DSE data quality.} If the DSE pull fails the go/no-go criteria in
    Section~\ref{sec:dse-data}, the target market changes to India: the NIFTY~50 index and large-cap NSE
    stocks, which are available through \texttt{yfinance} with more than 15 years of
    adjusted daily data. India is an emerging market, not a frontier one, so the transfer
    question becomes weaker, and the thesis would say so.
    \item \textbf{Floor-price periods.} About two and a half of the 13 or so DSEX years fall
    under floor prices. If excluding them leaves too little normal data, results are
    reported both with and without them.
    \item \textbf{Scarce labels.} Real anomaly labels on the DSE are few. The three label
    sources reduce dependence on any one of them, and injected anomalies give a controlled
    test.
    \item \textbf{Negative results.} If standardization or transfer does not help, the
    thesis reports it, along with the regimes or markets in which it helps or hurts.
    \item \textbf{Data licensing.} The DSE website terms have not yet been reviewed. Until
    they are, the code and download scripts will be released but raw DSE prices will not.
    \item \textbf{Compute.} The models are small enough to train on a laptop or on free
    Colab or Kaggle GPUs.
\end{itemize}

\section{Expected Contributions}
\begin{enumerate}
    \item A controlled test of GARCH standardization as the input to a deep anomaly
    detector.
    \item The first deep, out-of-sample anomaly detection study on the DSE among the work
    reviewed, with documented handling of floor prices, bonus shares and price limits.
    \item Evidence on transferring anomaly detectors from global markets to a frontier
    market.
    \item A measurement of how much point adjustment inflates performance on financial data.
    \item Public code and, if licensing allows, labelled DSE series.
\end{enumerate}

\section{Concluding Remarks}
The preliminary run already shows the problem the thesis addresses: on raw returns, a
statistical rule flags crisis years, while on standardized returns it flags days that were
unusual for their own period. The next step is the DSE data pull, followed by the labelling
rules, which will be fixed before any detector is trained.
